\documentclass[a4paper,11pt,openright,twoside]{report}      % Report class (dobre páxina, comezando capítulos na dereita)
%\parindent0pt \parskip10pt                                  % Para parágrafos tipo bloque

%%%%%%%%%%%% Cargamos os paquetes necesarios %%%%%%%%%%%%%%%%%%%%
\usepackage[utf8]{inputenc}				% para recoñecer vocais acentuadas e outros caracteres especiais
%\usepackage[bookmarksopen=true]{hyperref} \usepackage{bookmark}        % hiperligazóns no índice e outras utilidades
                                        % grande cantidade de erros ao empregar conxuntamente con \usepackage[galician]{babel} --- incluír só na versión final [!!]
\usepackage[nottoc]{tocbibind}          % para incluír a bibliografía no índice
\usepackage[galician]{babel}			% tradución de certos elementos tipográficos ao galego
\usepackage{graphicx}					% para incluír figuras
\usepackage{currfile}					% para usar o nome do arquivo en uso (versión)
%\usepackage[pass,showframe]{geometry}	% para mostrar as marxes do documento
%\usepackage{showlabels}                % para mostrar as etiquetas 
\usepackage{ccicons}                    % iconos das licencias Ceative Commons



%%%%%%%%%%%% Fornecer información sobre o documento %%%%%%%%%%%%%
\newcommand{\thesisTitle}{Exemplo de memoria para o STEMbach}                           % Título
\newcommand{\thesisName}{usando a \texttt{documentclass} de \LaTeX\ \texttt{report}}    % Autor/a
\newcommand{\thesisSubject}{Documentation}                                              % Temática
\newcommand{\thesisDate}{\today}                                                        % Data, \today para a data actual
\newcommand{\thesisKeywords}{kword1, kword2, kword3, kword4}                            % Palabras clave
\newcommand{\thesisVersion}{Primeiro borrador}                                          % Versión

% Define correctamente metadata para o documento PDF
\pdfinfo{                  
/Title (\thesisTitle)                       % dá erro se introducimos comandos 
/Author (X. O. Camanho) %\thesisName)       % ou caracteres especiais nos nomes
/Subject (\thesisSubject)
/Keywords (\thesisKeywords)
}


\title{ 
\vspace*{-4cm} {\small PROXECTO DE INVESTIGACIÓN}\\[3em]
\includegraphics[width=.6\textwidth]{ESCUDOROSALIA_1b.jpg}\\[1em] 
\bf \thesisTitle}				
\author{\thesisName}			        
\date{\thesisDate
\vfill 
\includegraphics[scale=.15]{STEMbach_logo.jpg}
}                           


%%%%%%%%%%%%%% Fin do preámbulo, comezo do documento %%%%%%%%%%%%%%

\begin{document}                        
\maketitle                              % Presenta a páxina de título
\pagenumbering{roman}                   % páxinas con numerais romanos para índices
\setcounter{page}{2}                    % facer que comece con "ii"

%%% información de licenza (CC 4.0) %%%%%%%%%%%
\newpage
\thispagestyle{empty}
\mbox{}
\vfill
\noindent\parbox{0.63\textwidth}{Documento preparado empregando o modelo para a realización da memoria de proxecto do {\bf Bacharelato de excelencia en ciencias e tecnoloxía (STEMbach)} do IES Plurilingüe Rosalía de Castro. \\[0.8em] {\Large\ccbync}\ \ {X. O. Camanho, 2020}
}

%%%%%%%%%%%%%

%%% páxina de información do documento %%%%%%%%%%%%
\newpage
%\thispagestyle{empty}
\begin{tabular}{p{0.4\textwidth}p{0.6\textwidth}}
\includegraphics[width=0.54\textwidth]{ESCUDOROSALIA_2.jpg} & 
\end{tabular}
\vskip1cm

\begin{tabular}{p{0.25\textwidth}|l}
%\includegraphics[scale=.1]{ESCUDOROSALIA_2.jpg} & \\
& \\
{\bf\hfill Título} & \thesisTitle\\[0.5em]
{\bf\hfill Autor} & \thesisName\\[0.5em]
{\bf\hfill Centro} & IES Plurilingüe Rosalía de Castro\\[0.5em]
{\bf\hfill Versión} & \thesisVersion, \today\\[0.5em]
{\bf\hfill Data} & \thesisDate\\[0.5em]
{\bf\hfill Titor} & \ldots\\[0.5em]
{\bf\hfill N\textordmasculine\ palabras} & \ldots\\[0.5em]
{\bf\hfill Licencia} & \ccbync\ CC-BY-NC 4.0 \\[0.5em]
& \\
\end{tabular}
%[2.5em]
\vfill
\begin{tabular}{p{0.25\textwidth} l}
\vskip-1.2cm\includegraphics[scale=.15]{STEMbach_big.jpg}
%logo_stembach_1.png} 
& \parbox{0.75\textwidth}{Esta memoria foi presentada\ldots Esta memoria foi presentada\ldots Esta memoria foi presentada\ldots Esta memoria foi presentada\ldots Esta memoria foi presentada\ldots Esta memoria foi presentada\ldots Esta memoria foi presentada\ldots Esta memoria foi presentada\ldots}
\end{tabular}




%%% Insert blank page %%%%%%%%%%%
%\newpage
%\thispagestyle{empty}
%\mbox{}
%%%%%%%%%%%%%



\tableofcontents                        % Mostra o índice 
\listoffigures                          % Mostra o índice de figuras
\listoftables                           % Mostra o índice de táboas


\chapter*{Resumo}						% Resumo do plantexamento inicial e conclusións

Bla bla bla\ldots

\section*{\Huge Abstract}\vskip.8cm				% en inglés
%\chapter*{Abstract}


Bla bla bla\ldots

\chapter{Encabezamento principal}       % Mostra un encabezamento de "capítulo"
\pagenumbering{arabic}                  % Numeración indo-arábiga comezando en 1

\begin{flushright}                      % Aliña o texto na marxe dereita
    \textit{\small                      % Texto en cursiva máis pequeno
    ``We are like dwarfs sitting on the shoulders of giants.\\ We see more, and things that are more distant,\\ than they did, not because our sight is superior\\ or because we are taller than they, but because\\ they raise us up, and by their great stature add to ours.''
    }\vskip2mm
    
    John of Salisbury
\end{flushright}

A maioría do que segue aplica tanto ás clases \texttt{article} e \texttt{book} 
como a clase \texttt{report}. Porén, na clase \texttt{article}, a posición por defecto da información de título é ao comezo da primeira páxina de texto e non nunha páxina separada. Por outra banda, non é habitual incluír un índice coa clase \texttt{article}.

Bla bla bla
 
\section{Encabezamento de sección}      % Mostra un encabezamento de "sección"
Están disponibles os seguintes comandos de sección:

 \begin{verbatim}                                         
        \part 
        \chapter                            
        \section 
        \subsection  
        \subsubsection  
        \paragraph 
        \subparagraph 
 \end{verbatim}                         % entre  \begin{verbatim}   ...    \end{verbatim}    
                                        % mostra todo sexa texto ou comandos             

Fágase notar que, a diferencia das clases \texttt{book} e \texttt{report}, a clase
\texttt{article} non ten comando \verb+\chapter+.

Un dos xeitos de incluír citas intercaladas no texto é empregando o entorno \verb+quote+

\begin{center}
\verb+\begin{quote}+

$\ldots$

\verb+\end{quote}+
\end{center}

\begin{quote}                           % O seguinte texto vai estar separado 
                                        % e indentado
                                        % \\ provoca o comezo dunha nova liña
    We are like dwarfs sitting on the shoulders of giants. We see more, and things that are more distant, than they did, not because our sight is superior or because we are taller than they, but because they raise us up, and by their great stature add to ours.
    
    \hfill John of Salisbury
\end{quote}                             % Fin do texto indentado

Outro comando útil é \verb+\footnote{}+ para incluír notas\footnote{Coma esta} a pé de páxina.


No que respecta ás fontes, estas pódense presentar en \textit{cursiva} ou \textbf{negrita} e moduladas en tamaño: {\tiny \verb+\tiny+}, {\scriptsize \verb+\scriptsize+}, {\footnotesize \verb+\footnotesize+}, {\small \verb+\small+}, {\normalsize \verb+\normalsize+}, {\large \verb+\large+}, {\Large \verb+\Large+}, {\LARGE \verb+\LARGE+}, {\huge \verb+\huge+}, {\Huge \verb+\Huge+}

 
\subsection{Encabezamento de subsección} 

Bla bla bla

\subsection{Exemplos}

Esta linguaxe de composición de textos é especialmente útil cando temos que incluír fórmulas no texto que poden inserirse entre \verb+$...$+ no medio dunha liña, por exemplo, $E=m\cdot c^2$; ou entre \verb+$$...$$+ como fórmula indentada:

$$E=m\cdot c^2$$

Para fórmulas numeradas empregaremos mellor o entorno \verb+equation+:

\begin{equation}
\label{mc2}
    E=m\cdot c^2
\end{equation}
o que nos permite ademais etiquetala e referirnos a esta fórmula noutros lugares do texto (ver ecuación \ref{mc2}). Do mesmo xeito podemos etiquetar capítulos, seccións, etc. para facer referencias cruzadas ao longo do texto (e.g. ver capítulo \ref{metodos}).




\begin{table}[ht]
    \centering
    \begin{tabular}{c|c}
        A & B \\
    \hline
        ss & ad \\
        as & dd
    \end{tabular}
    \caption{Este é o pé da táboa}
    \label{tab:exemplo}
\end{table}





Citas bibliográficas...





Ver \cite{einstein}

\begin{figure}[ht]
    \centering
    \includegraphics[scale=.15]{STEMbach_big.jpg}
    \caption{Este é o pé da imaxe}
    \label{fig:exemplo}
\end{figure}


%%%%%%%%%%%%%%%%%%%%%%%%%%%%%%%%%%%%%%%%%%%%%%%%%%%%%%
% Os títulos dos seguintes capítulos 
% están tomados dos requisitos establecidos 
% na resolución que regula o STEMbach.
% Poderíanse tamén estruturar como seccións 
% (sen capítulos) para documentos máis breves
%
%%%%%%%%%%%%%%%%%% INTRODUCIÓN %%%%%%%%%%%%%%%%%%%%%%%
\chapter{Introdución}
\label{intro}


\chapter{Contextualiación e xustificación} % da investigación realizada
\label{contexto}





\chapter{Materiais e métodos}
\label{metodos}


\chapter{Resultados e conclusións} % (incluindo unha reflexión sobre o proceso de investigación)
\label{conclusions}


\chapter{Proposta de continuidade}
\label{continuidade}





%%%%%%%%%%%%%%%%%% BIBLIOGRAFÍA %%%%%%%%%%%%%%%%%%%%%%%%
\begin{thebibliography}{9}
\bibitem{latexcompanion} 
Michel Goossens, Frank Mittelbach, and Alexander Samarin. 
\textit{The \LaTeX\ Companion}. 
Addison-Wesley, Reading, Massachusetts, 1993.
 
\bibitem{einstein} 
Albert Einstein. 
\textit{Zur Elektrodynamik bewegter K{\"o}rper}. (German) 
[\textit{On the electrodynamics of moving bodies}]. 
Annalen der Physik, 322(10):891–921, 1905.
 
\bibitem{knuthwebsite} 
Knuth: Computers and Typesetting,
\\\texttt{http://www-cs-faculty.stanford.edu/\~{}uno/abcde.html}
\end{thebibliography}
 
 
\end{document}                          % última liña requerida para rematar o documento 
                                        % (o que segue non é tido en conta para a 
                                        % elaboración do documento)
